\scanpage{5}
\mspage{3}
Soient $x,y\in L$, on dit que $y$ est une \uline{spécialisation} de $x$ si \deleted{$\unsure{\exists}$}

a) $y$ est singulier, et b) $\exists I\in S$ tel que $x\in\tilde I^\circ$, $y\in\delta I$.

\deleted{
\uline{Axiomes}\deleted{\unsure{s}}

$F_1)$ Soit $I\in S$, alors \deleted{$\unsure{\exists}$} $\tilde I^\circ$ ($\underset{\text{déf}}{=}\tilde I\setminus\delta I$) est non vide, et si $z\in\tilde I^\circ$, \deleted{$\unsure{\exists}$} $\delta I=\{x,y\}$,
\[
\exists!\ I_{x,z}\in S,\ I_{y,z}\in S\quad\text{tels que}
\]
\[
\tilde I_{x,z},\tilde I_{y,z}\unsure{\in|\subset}\tilde I,\qquad
\delta I_{x,z}=\{x,z\},\quad\delta I_{y,z}=\{y,z\}.
\]
De plus, on a alors
\[
\tilde I_{x,z}\cap\tilde I_{y,z}=\{z\}
\]
}
\margin{décomposition des segments

\uuline{décomposition} des segments

Mais attention, je ne dis pas que $\tilde I=\tilde I_{x,z}\cup\tilde I_{y,z}$ !!!}
\margin{Ajouter: $\forall\unsure{\tilde J|J}\in S$ tel que $\tilde J\subset\tilde I$, $\delta\tilde J\ni z$, on a $\tilde J\subset\tilde I_{x,z}$ ou $\tilde J\subset\tilde I_{y,z}$. \deleted{\unsure{Alors}} \deleted{\unsure{il}}}

\deleted{
$F_2)$ \deleted{Soient $x,y$ lieux réguliers. Alors, \unsure{il} existe $I\in S$ tel que $x,y\in\tilde I^\circ$, alors}
\inserted{Soient $I\in S$, et \deleted{$\unsure{x,y}$} $\unsure{\Phi|\varphi}\in\tilde I^\circ$}
$\exists$ \deleted{$\unsure{I\in S}$} $I'\in S$ tel que
a) \deleted{$\unsure{x}$} $\unsure{\Phi|\varphi}\in\tilde I'^\circ$,
b) $\tilde I'\subset\tilde I^\circ$.

\inserted{\uline{Corollaire 1} $\tilde I$ est \uline{infini}.

\uline{Corollaire 2} Si $S\neq\varnothing$, $S$ est \uline{infini}\ldots
(car l'ens des $\unsure{\tilde J|J}\in S$ tels que $\tilde J\subset\tilde I$ est infini\ldots)}

\deleted{$\unsure{F_2}$} $F_2)$ (Réciproque de $F_1$) Soient $I',I''\in S$, tels que
\[
\tilde I'\cap\tilde I''=\deleted{\unsure{\{z\}}}\delta I'\cap\delta I''=\{z\},
\]
avec $z$ \uline{régulier}. Alors $\exists!\ I\in S$, tel que
\[
\tilde I'\cup\tilde I''\subset\tilde I\deleted{\unsure{*}}
\quad\text{et}\quad
\delta I=(\delta I'\cup\delta I'')\setminus\{z\},
\]
et on \inserted{a, si $\delta I'=\{x,z\}$, $\delta I''\deleted{\unsure{=}}\ =\{y,z\}$, $\unsure{I'|\tilde I'}=\tilde I_{x,z}$, $\unsure{I''|\tilde I''}=\tilde I_{y,z}$.}

\deleted{$F_3)$ Soit $x\in\unsure{L}$}
\inserted{\unsure{cf plus bas}}

$F_3)$ Soit $I\in S$, $\Phi$ partie finie de $\tilde I^\circ$,
Alors $\exists I'\in S$ tel que $\tilde I^\circ\supset\tilde I'\supset\Phi$.
}
\margin{\uuline{recollement} des segments}
\margin{\uuline{\unsure{Filtration}}}
\margin{\deleted{\unsure{l'ens des} $\tilde I$ \unsure{tels que} $x\in\tilde I^\circ$ \unsure{est filtrant}, si on a $x\in\tilde I'^\circ\cap\tilde I''^\circ$,
$\exists\tilde I$ tel que
$x\in\tilde I^\circ\subset\tilde I'^\circ\cap\tilde I''^\circ$.
Il en résulte qu'il \unsure{existe} $I$, \unsure{avec}
\deleted{$\unsure{\tilde I\subset\tilde I'^\circ}$}
$\tilde I\subset\tilde I'^\circ\cap\tilde I''^\circ$,
$x\in\tilde I^\circ$.}}
\mspage{4}

\rotatebox{90}{\unsure{EXTRA STRONG}}
\rotatebox{90}{\unsure{EXTRA STRONG}}
